\section{Related Work}

xView3 frames dark-vessel detection as localization in Sentinel-1 SAR rather
than box classification~\cite{paolo2022xview3}. A CenterNet-style heatmap fits
the point annotations and distance metric without inventing box extents
~\cite{zhou2019objects}. We use this established detector form to hold the
task head fixed.

The ViT track compares supervised ImageNet AugReg~\cite{steiner2021trainvit},
optical SatDINO trained on fMoW-RGB~\cite{straka2025satdino}, and SARMAE trained
on SAR-1M~\cite{liu2026sarmae}. SatDINO and SARMAE also differ in objective, so
their gap is a released-checkpoint comparison rather than a pure modality
intervention. The CNN track uses ConvNeXt-V2-Base~\cite{woo2023convnextv2}.
Its optical and SAR arms come from matching BigEarthNet Sentinel-2 and
Sentinel-1 land-cover models~\cite{sumbul2019bigearthnet,clasen2025reben}; this
pair gives the cleaner modality contrast. ImageNet controls match the generic
source dataset but not full training history across architectures.
